% Generated draft: requires mathematical and source review.
\documentclass[11pt]{article}
\usepackage[margin=25mm]{geometry}
\usepackage{fontspec}
\setmainfont{Latin Modern Roman}
\usepackage{amsmath,amssymb,mathtools,unicode-math}
\setmathfont{Latin Modern Math}
\usepackage{enumitem,needspace,xurl,hyperref}
\setlength{\parindent}{0pt}
\setlength{\parskip}{5pt}
\setlength{\emergencystretch}{3em}
\newcommand{\N}{\mathbb{N}}
\newcommand{\Z}{\mathbb{Z}}
\newcommand{\Q}{\mathbb{Q}}
\newcommand{\R}{\mathbb{R}}
\newcommand{\C}{\mathbb{C}}
\newcommand{\F}{\mathbb{F}}
\newcommand{\A}{\mathbb{A}}
\newcommand{\PP}{\mathbb{P}}
\newcommand{\E}{\mathbb{E}}
\newcommand{\eps}{\varepsilon}
\newcommand{\dd}{\,\mathrm{d}}
\begin{document}
\section*{UnsolvedMath ID: 30001728}
\textbf{Source problem code:} catalogue entry
\subsection{Definitions and mathematical statement}
\textbf{Research question.} What can examples of interpolation in late-medieval double-entry astronomical tables reveal about (i) the algorithms that shape tables, (ii) the treatment of discrete data as continuous, and (iii) the mathematical and social uses of tables?

\textbf{Scope and unresolved definitions.} This is a historiographical research question. The supplied record does not identify a corpus of tables, particular interpolation procedures, the meaning of ``double-entry'' in the intended examples, or precise chronological and geographical boundaries for ``late-medieval.'' It also does not specify how algorithms are to be reconstructed, what constitutes treating discrete data as continuous, or what evidence would establish the mathematical and social uses of the tables. These terms therefore remain descriptive rather than formally defined here.

\textbf{Limits of mathematical formulation.} No mathematical objects with specified domains, hypotheses, boundary conditions, quantifiers, or conclusion are supplied from which a definite theorem or conjecture can be stated. The question is retained as a research question; a precise mathematical proposition cannot be extracted without additional source evidence and an explicit intended target.

The registry labels the item ``open.'' Its supplied literature assessment describes the extract as a fragment of historical discussion without a complete claim having a known proof or disproof, and retains the open label under its treatment of unclear extracts. Separately, the supplied statement-verification field labels the wording \texttt{corrected\_verified} and reports that the official report confirms the historiographical question and that an extraction label was removed. These are supplied curation assertions, not independent verification or a proof-based classification.
\subsection{Short English statement}
The question asks what interpolation in late-medieval double-entry astronomical tables can teach us about table-making algorithms, the interpretation of discrete data as continuous, and the mathematical and social roles of tables. The supplied material supports a historical research question but does not specify a mathematical conjecture. Its scope, terminology, and evidentiary criteria require human review.
\subsection{Sources}
\begin{itemize}
\item Local UnsolvedMath registry, permanent ID \texttt{30001728}, title \emph{Interpolation in Medieval Astronomical Tables}, source problem code \texttt{catalogue entry}, in the collection \emph{Oberwolfach Reports Open Problems}. The record gives proposed year 2011 and names no proposer.
\item Supplied source citation: \emph{Mini-Workshop: History of Numerical and Graphical Tables} (2011), \emph{Oberwolfach Reports}, pp.~639--689, DOI \texttt{10.4171/owr/2011/12}. Supplied source URL and literature-source link, labeled ``Oberwolfach Report 12/2011'': \url{https://doi.org/10.4171/owr/2011/12}. The supplied background describes Oberwolfach Reports as documenting mathematical workshops and their open-problem discussions.
\item Literature assessment dated exactly \texttt{2026-08-21}: the supplied curation labels the item open because the historical extract does not supply a complete claim with a known proof or disproof. This is supplied dated curation, not independently verified. The background expressly requests independent verification of time-sensitive status and cited claims.
\item The supplied statement-verification note reports confirmation of the historiographical question by the official report. No independent source visit or verification is asserted in this draft.
\end{itemize}
\end{document}
